\daymark{125}{Reed 6, Sixthday, Y.H. 813}

The morning came up off \place{the Ilm} with mud in it and cut grass, the same thin mineral edge the river left on everything it touched, and \person{Sella Vail} carried \person[Corin]{the boy} down the road to \place{the Weir} at the half-pace she kept from every open slate. He was one hundred and twenty-five days old, four months and five days, and nobody who carried him had any use for the number. The reed work had come downriver in the night the way it always came, not as a figure but as a direction: bundles going out and bundles coming back, judged by the arm that brought them and by nothing else. At the edge of the work she set him where no slate lay, gave the instruction to the wool and the river wind and never to the basket, and went a half-pace away to stand where she could see the boards and the posts and the men upon them. The basket sat level. Nothing in it slopped. His hand outside the wool went up an inch toward the low mutter of the water and came down, because two inches was the whole of his reach and the water was a league away.

\place{the reed beds} went on behind the yard in long grey ranks, green at the broken ends where the knife had been and pale at the frost-bitten tips, the colour of old bone laid out in ruts. The traps stood in the gaps of the weir with their mouths upriver, and \place{the plank walk} went out over the water on posts with salt dried white where the flood had touched them, a dry, honest line of a thing that would bear a man's weight until the day it would not and would give no notice first. A girl with a creel went along it ahead of the others, her step quick and her eyes on the boards. The boy watched the boards as a man watches a floor he has not yet trusted, but he had no hand to put on one and no voice to say which would hold.

A reed-cutter came along the planks with a bundle on his shoulder, the broken ends wet and dark against his coat, and set the bundle down on the flat stone where the hearth-maids had stood in Ilmen to bless the traps. He took a short length of wool or reed-fibre from inside his coat, gone grey with handling, and straightened it on his thumb the way a man straightens a thing he means to trust. The knots stood at the lengths a man's arm could judge. He turned the cord once so they stood in the small light off the water, and the boy saw the shape of them without reading them, which was the whole of what the cord was for. The reed-cutter tied one low down with his teeth and a pull of the thumb, bit the end level, and put the cord back inside his coat. He did not show it to anyone and he did not count it where anyone could hear him count.

A weir woman with salt dried white at her wrists came up the boards and asked how many had gone up yesterday.

“You can ask the cord,” the reed-cutter said.

“I can't read a cord.”

“No more can the Court.” He took up his bundle and went to lift a trap, and the question went along with him and down into the water and did not come up again with the eels.

No figure came out of it. No knot was counted and no name was given, and the boy did not lay the cord across the place in \place{the worked strips} where a hand had gone in with the seed and had not come out with the rest. A man's tally of reeds was a man's tally of reeds; an absent hand had no figure even in a cord that no Court thumb would read. He saw the cord as he had seen a chart closed with a neat line over a note nobody had opened --- exact, complete, worth nothing to anyone and most of all to him, who could not have closed two inches on it or told a soul what it kept --- and set it on a place of its own.

It was the third such place he kept. A woman at the eel boards had a grey wool cord at her waist and a second twist twisted twice beneath it, and when a figure came near her thumb went under her apron to the first and stayed longer than the figure required; neither knot had ever been counted by any thumb but her own, and neither cord had ever been carried upriver to \place{Nine Pots} or lifted at a door. Another woman's knot-string lay on a board beside hers and did not touch it. Six pages in a little tin, written down in a second hand, lay turned writing-down toward a floorboard in a house he could not enter, and a round pewter tin above a bed nobody used kept its pale fold and its thumb-bright lid shut through more mornings than anybody kept count of. No tin was opened. No box was opened. Every record in the district was kept in a material no Court thumb could spell, and none of them touched another.

\scenebreak

\person{Orra Weir} stood at the lower boards with her board on its side and a cord at her waist, pulling one knot down with her thumb and setting it among the others. Her hands were like old rope and the pale eel scars lay uncovered across both forearms, and the grey streak showed first through the black hair when she turned her head, the way it had shown first at the blessing and every morning since. The nasal half-song of the Fen came into her voice where she was tired, and she had never entirely lost it. The weir woman with the salt-white wrists asked what the levy would take now the hands were out of the reeds.

“The ninth of Ilmen,” Orra said. “The collector's boat came round the bend on the flood and had six oars. One man stood at the stern and never touched them.”

“That isn't the levy.”

“It's a date.” She pulled another knot down. “The river takes what the Court doesn't. You can shout at either one. I've tried.”

A date given where a figure was asked for was an answer with the question taken out of it, a chart with the finding left off and the neat margin left in; the boy had learned to read the shape in a woman's face before he could read the words on a slate, and he set her saying where it belonged and made no cause of it. Her thumb went to the cord again as the sentence ended, and whether that was a lie or only the habit of counting he could not tell, and he did not make a cause of that either.

When the reed-cutter lifted his bundle off the stone she watched what there was to watch at the edge of the shade and got no answer from the basket. Nobody asked her and she did not ask. The shape of the last knot on her cord was legible without its name --- double turn, tail tucked, left over right, the small extra weight of a thumb --- and the boy did not spend it. The word forty stayed where it had stayed since the year of the flower, unspoken in a morning that had a great many figures in it and not one he could enter. Her board lay beside the eel board and did not touch it.

\person{Mira Silt} worked the eels at the outer boards with pale slime dried in long grey strokes on her sleeves and river silt dried pale in the hollow of her neck, her head held a little to the right where the top of her left ear had been, her small smile kept on her face the way another woman keeps a door on the latch. A weir woman asked whether the run would be thin now that the hands were in the reeds.

“Thin? Who's to say thin?” Mira said. “It comes when it comes. Ask the river.”

She flicked a string of eel-bones against her thumb while deciding whether to lie, moved two of them from hand to hand, and put one away.

A buyer in a coat with the river still in its hem named a figure for two pair. Mira's thumb went under her apron to the grey cord at her waist and stayed there longer than the figure required, and when she shifted a creel to reach into a basket the second cord twisted twice beneath it showed, and neither knot was counted and the knot for the thin child upriver was not touched. “Two marks the pair and two short,” she said. “You've said it twice and I've counted it twice, and I'll count it again after I've eaten. There's nothing in the sentence that wants an answer.” She did not look into the basket. Her not-looking at the boards was a keeping, set on a place of its own and joined to no pricing; her board lay beside Orra's and did not touch it.

Beyond \place{the Cross} the stalls had come up in the ordinary way, reeds laid over eels in one long noisy row where a buyer could not tell which hand had cut what, and the market sold over the top of a man with a slate as though a man with a slate were a thing every landing had. Below \place{the north bridge} at \place{the Fishstairs}, on \place{the out-of-line landing} where a step should have been, the soft-cuffed clerk was still set. His cuff was grey from washing and grey again by midday. He held his slate to his forearm for a measure nobody had called, and his hand went to the slate before his eyes went to the stalls, every time, so that the reaching came before the looking and the looking never changed the slate. Nobody stepped around him. Nobody asked what he was set there for.

A girl with a creel went over the landing sideways, caught the rail with the hand that was not carrying, and went on selling. Nobody made a fall of it. The boy put the order of the clerk's hand on its own place, apart from the women's hands at the boards and apart from the man who would come up the road with a clean thumb, and wrote no cause across any of them. A measure set down before anybody called it was a measure set down ahead of something; he had kept that as a shape for a hundred and twenty-five days and had never made a rule of it. The market's noise went over the slate in a sheet, the way water went over a stone and left it where it was.

\scenebreak

\person{Berr Clay} came up short off \place{the clay road} a little after midday with the red wool cord of his office folded under his coat and the green of tasted glaze at one corner of his mouth, hot clay over the smell of nothing at all, the most expensive smell in \place{Grath}. His hands were large and clean. He wiped his thumb twice on his apron before touching nothing and held a narrow slate unopened against his forearm, passed a finger over the figure and let his eyes rest on the writing a breath longer than the writing required, the way a man checks a dish he has set down to be sure it will not slop.

“A number goes into a slate, not into a bundle,” he said. “If your house wants it again, there'll be a slate at the watch.”

He read one in twelve, as the notice says, low into the two hundred, called the levy and said it was not called by him. A man brought a number up out of a fist and asked what the reeds would cost now the hands were out of them. Berr told him there would be a slate at the watch --- \emph{your house}, never \emph{I} --- and took himself out of his own sentence as neatly as a surgeon takes a hand out of a field before the dressing goes on. The number went from the fist to his mouth to a shoulder going up the road and arrived nowhere. He did not look into the basket. He went off up the clay road toward \place{High Kiln} with the box in his house still locked and the true subtraction inside, believing he had done the district a kindness.

A figure read aloud was not the figure paid. A number put into a slate was not a number put into a strip. The boy kept the three things on three places --- the cord in the man's coat, the tally in the slate, the figure that went nowhere --- and no hand of his could have laid one across another if his life had depended on it, which it did, though nobody under any roof in Sill would have known what the dependence was.

The count stood at one hundred and twenty-five where one hundred and twenty-four had stood the day before, exact, unspent and spoken by nobody. No loaf passed. No slate was opened within reach of the basket. No sixth sighting of the neat dark hand came; the foundling was nowhere at the market's edge, \place{Sill} held no collector anybody could have mistaken for a market woman, and the woman more mornings dead than any living person keeps count of was not named by any person in the morning. The reed-cutter took up his bundle and went along the boards, the clerk's hand went to the slate before his eyes went to the stalls, and the market sold its reeds and its eels over the top of him as though a man with a slate were a thing a landing always had. The cord went back inside the man's coat with its knots unread, and nobody stepped around him.
